% from ega4-11-fr.tex
\oldpage[IV-3]{118}

En vertu du théorème de platitude générique (6.9.1) appliqué à l'anneau
intègre $A/\mathfrak p$, à la $(A/\mathfrak p)$-algèbre de type fini
$B/\mathfrak pB$ et au $(B/\mathfrak pB)$-module de type fini
$M/\mathfrak pM$, il existe un $h\in A-\mathfrak p$ tel que, si $\bar h$ est
son image canonique dans $A/\mathfrak p$,
$(M/\mathfrak pM)_{\bar h}$ soit un $(A/\mathfrak p)$-module plat. D'autre
part, comme $M_{\mathfrak q}$ est un $A_{\mathfrak p}$-module plat, et par
suite un $A$-module plat ($\mathbf 0_{\mathrm I}$, 6.3.1), on a
$\operatorname{Tor}_1^A(M_{\mathfrak q},A/\mathfrak p)=0$, ce qui s'écrit
aussi
$(\operatorname{Tor}_1^A(M,A/\mathfrak p))_{\mathfrak q}=0$. Mais comme $A$
et $B$ sont noethériens, $\operatorname{Tor}_1^A(M,A/\mathfrak p)$ est un
$B$-module de type fini, donc ($\mathbf 0_{\mathrm I}$, 5.2.2) il existe
$g\in B-\mathfrak q$ tel que
$(\operatorname{Tor}_1^A(M,A/\mathfrak p))_g=0$. D'ailleurs, on a
$(M/\mathfrak pM)_{\bar h}=(M/\mathfrak pM)_h$ ($M/\mathfrak pM$ étant
considéré dans le second membre comme un $A$-module)~; de plus,
$(M/\mathfrak pM)_h$ étant un $B$-module,
$(M/\mathfrak pM)_{hg}$ est encore un $(A/\mathfrak p)$-module plat, car il
s'écrit $(M/\mathfrak pM)_{\overline{hg}}$, où $\overline{hg}$ est l'image
canonique de $hg$ dans $B/\mathfrak pB$, et il suffit d'appliquer
($\mathbf 0_{\mathrm I}$, 6.3.2)~; enfin, on a
$(\operatorname{Tor}_1^A(M,A/\mathfrak p))_{hg}=0$ et
$hg\in B-\mathfrak q$ puisque $h\in A-\mathfrak p$. C.Q.F.D.

\begin{corollary}[11.1.2]
\label{IV.11.1.2-fr}
Soient $Y$ un préschéma localement noethérien, $f:X\to Y$ un morphisme
localement de type fini, $\mathcal F$, $\mathcal F'$ deux
$\mathcal O_X$-Modules cohérents,
$u:\mathcal F'\to\mathcal F$ un homomorphisme de $\mathcal O_X$-Modules. On
suppose que $\mathcal F$ est $f$-plat. Alors l'ensemble $U$ des $x\in X$ tels
que, en posant $y=f(x)$, l'homomorphisme
\[
  u_x\otimes1:\mathcal F'_x\otimes_{\mathcal O_y}k(y)\longrightarrow
  \mathcal F_x\otimes_{\mathcal O_y}k(y)
\]
soit injectif, est ouvert dans $X$.
\end{corollary}

En effet, soit $\mathcal N$ (resp. $\mathcal P$) le noyau (resp. le conoyau) de
$u$~; appliquons ($\mathbf 0_{\mathrm{III}}$, 10.2.4)\footnote{Voici une
démonstration de ($\mathbf 0_{\mathrm{III}}$, 10.2.4), qui n'a pas été publiée
dans l'\emph{Algèbre commutative} de N.~Bourbaki. Compte tenu de
($\mathbf 0_{\mathrm I}$, 6.1.2), il suffit de voir que b) entraîne a). Posons
$P=\operatorname{Im}(u)$, $Q=\operatorname{Coker}(u)$,
$R=\operatorname{Ker}(u)$. Le composé
$M\otimes k\to P\otimes k\to N\otimes k$ est injectif, et
$M\otimes k\to P\otimes k$ est surjectif, donc
$P\otimes k\to N\otimes k$ est injectif et
$M\otimes k\to P\otimes k$ est bijectif. La suite exacte
$0\to P\to N\to Q\to0$ donne la suite exacte
\[
0=\operatorname{Tor}_1^A(N,k)\to\operatorname{Tor}_1^A(Q,k)
\to P\otimes k\to N\otimes k\to Q\otimes k\to0
\]
et puisque $P\otimes k\to N\otimes k$ est injectif, on a
$\operatorname{Tor}_1^A(Q,k)=0$~; comme $Q$ est un $B$-module de type fini,
($\mathbf 0_{\mathrm{III}}$, 10.2.2) montre que $Q$ est un $A$-module plat~;
alors $P$ est aussi un $A$-module plat par ($\mathbf 0_{\mathrm I}$, 6.1.2).
La suite $0\to R\to M\to P\to0$ étant exacte, il en est de même de
$0\to R\otimes k\to M\otimes k\to P\otimes k\to0$ par
($\mathbf 0_{\mathrm I}$, 6.1.2)~; puisque $M\otimes k\to P\otimes k$ est
bijectif, on a $R\otimes k=0$~; mais comme $B$ est noethérien, $R$ est un
$B$-module de type fini, donc on a $R=0$ en vertu du lemme de Nakayama.} aux
anneaux locaux $\mathcal O_y$ et $\mathcal O_x$ et aux
$\mathcal O_x$-modules $\mathcal F'_x$ et $\mathcal F_x$~: dire que
$u_x\otimes1$ est injectif équivaut à dire que $\mathcal N_x=0$ et que
$\mathcal P$ est $f$-plat au point $x$. Or comme $\mathcal N$ et $\mathcal P$
sont cohérents ($\mathbf 0_{\mathrm I}$, 5.3.4), l'ensemble des $x$ où
$\mathcal N_x=0$ est ouvert ($\mathbf 0_{\mathrm I}$, 5.2.2), et l'ensemble
des $x$ où $\mathcal P$ est $f$-plat est ouvert par (11.1.1)~; d'où la
conclusion.

En particulier~:

\begin{corollary}[11.1.3]
\label{IV.11.1.3-fr}
Soient $Y$ un préschéma localement noethérien, $f:X\to Y$ un morphisme plat et
localement de type fini, $g$ une section de $\mathcal O_X$ au-dessus de $X$.
L'ensemble des $x\in X$ tels que $g_x\otimes1$ soit non diviseur de zéro dans
$\mathcal O_x\otimes_{\mathcal O_{f(x)}}k(f(x))$ est ouvert dans $X$.
\end{corollary}

Il suffit d'appliquer (11.1.2) à l'endomorphisme de $\mathcal O_X$ défini par
$g$.

\begin{corollary}[11.1.4]
\label{IV.11.1.4-fr}
Soient $Y$ un préschéma localement noethérien, $f:X\to Y$ un morphisme
localement de type fini, $\mathcal F$ un $\mathcal O_X$-Module cohérent et
$f$-plat. Soit $(g_i)_{1\leq i\leq n}$ une suite de sections de $\mathcal O_X$
au-dessus de $X$. Alors l'ensemble $U$ des $x\in X$ tels que la suite
$((g_i)_x\otimes1)$ soit
$(\mathcal F_x\otimes_{\mathcal O_{f(x)}}k(f(x)))$-régulière est ouvert dans
$X$.
\end{corollary}

